\chapter{Etude de la problematique et objectifs}

\section{Problematique}
La migration d'une VM classique vers OpenShift Virtualization est un processus sensible qui depend de plusieurs facteurs: format de disque, ressources CPU/RAM, connectivite reseau, stockage disponible dans le cluster et strategie de migration.

Dans une approche manuelle, ces verifications sont longues et sujettes aux erreurs. Il devient donc necessaire d'introduire une solution capable de standardiser la preparation, d'anticiper les blocages et d'aider a la prise de decision.

\section{Objectif general}
L'objectif general est de developper une plateforme de migration intelligente de VMs vers OpenShift, avec un pipeline complet allant de la decouverte a la migration.

\section{Objectifs specifiques}
Les objectifs specifiques du projet sont:
\begin{enumerate}[leftmargin=1.2cm]
  \item decouvrir les VMs sources (VMware Workstation, KVM/libvirt),
  \item analyser la compatibilite avec OpenShift Virtualization,
  \item produire un plan de conversion technique,
  \item recommander une strategie de migration,
  \item executer et suivre la migration via les outils OpenShift/KubeVirt.
\end{enumerate}

\section{Demarche adoptee}
La demarche suivie repose sur une progression iterative:
\begin{itemize}[leftmargin=1.2cm]
  \item stabiliser les fondations backend/frontend,
  \item integrer les briques metier une a une,
  \item tester sur des VMs reelles de taille differente,
  \item corriger les blocages identifies sur le terrain.
\end{itemize}
